\section{Dataset}
\label{sec:dataset}

\subsection{Source Cohort and Annotation Protocol}
\label{sec:dataset_source}

We use VinDr-Mammo version 1.0.0, a credentialed open-access collection of
full-field digital mammography distributed through PhysioNet
\cite{nguyen2023vindr,pham2022vindrphysionet}.  The release contains 5,000
examinations and 20,000 DICOM images.  Each examination comprises the standard
left and right craniocaudal (CC) and mediolateral oblique (MLO) views.  Images
were acquired retrospectively at Hospital 108 and Hanoi Medical University
Hospital in Vietnam between 2018 and 2020.  Random sampling from the two
hospitals' archives yielded a mixture of screening and diagnostic examinations
acquired on Siemens, IMS, and Planmed systems.  Source image heights and widths
vary substantially, with dataset-wide maxima of 3,580 and 2,812 pixels,
respectively.  Thus, the common practice of forcing a mammogram into a
$640\!\times\!640$ square would retain only about one-fifth of the samples
along each axis and change the image aspect ratio before feature extraction.

Every examination was independently interpreted by two radiologists with at
least ten years of experience; disagreements were resolved by a third reader.
The release provides breast-level density and BI-RADS assessment as well as
lesion-level bounding boxes.  Finding labels include mass, calcification,
asymmetry subtypes, architectural distortion, suspicious lymph node, and
associated skin or nipple changes.  Boxes were drawn for findings assessed as
BI-RADS 3--5; BI-RADS 2 findings were intentionally not localized.  The
annotation file contains 1,226 mass boxes on 1,113 distinct mammograms.  We use
these boxes as radiologist-defined localization targets, not as
pathology-confirmed cancer labels: the release does not include biopsy results
and should not be interpreted as a direct cancer-diagnosis endpoint
\cite{nguyen2023vindr}.

\subsection{Task Definition and Partitions}
\label{sec:dataset_task}

Our endpoint is class-agnostic mass localization.  We retain a finding when
\texttt{Mass} occurs in its released category list and map every retained box
to a single detector class.  A mammogram without a retained mass is therefore
mass-negative even when it contains another annotated finding.  Unlike the
positive-heavy subsets reviewed in Sec.~\ref{sec:background_vindr_models}, we
retain all 20,000 mammograms.  This preserves the original prevalence of normal
and non-mass images and makes false-positive performance part of the task.

VinDr-Mammo supplies an official examination-level partition of 4,000 training
examinations and 1,000 test examinations, created to preserve breast-level
assessment, density, and finding distributions \cite{nguyen2023vindr}.  We
leave its 1,000-examination test set untouched.  For checkpoint selection, we
reserve 15\% of the official training examinations as validation using seed
123.  Sampling is stratified by the sorted set of breast-level BI-RADS values
present among the four views of each examination; all four views then inherit
the same assignment.  The resulting identifiers are persisted and reused for
every resolution.  This protocol connects directly to the comparability issue
identified in Sec.~\ref{sec:background_vindr_models}: neither correlated views,
negative-case policy, nor test membership changes between experiments.

\begin{table}[t]
\centering
\caption{Analysis cohort after selecting Mass annotations. A mass-negative
image may contain a different VinDr finding.}
\label{tab:dataset_split_counts}
\scriptsize
\setlength{\tabcolsep}{3.2pt}
\begin{tabular}{@{}lrrrrr@{}}
\toprule
Split & Exams & Images & Mass-neg. & Mass-pos. & Boxes \\
\midrule
Train & 3,400 & 13,600 & 12,857 & 743 & 829 \\
Validation & 600 & 2,400 & 2,249 & 151 & 160 \\
Test & 1,000 & 4,000 & 3,781 & 219 & 237 \\
\midrule
Total & 5,000 & 20,000 & 18,887 & 1,113 & 1,226 \\
\bottomrule
\end{tabular}
\end{table}

\subsection{DICOM Conversion and Breast Standardization}
\label{sec:dataset_preprocessing}

We materialize the cohort with our Default Research Dataset v2 pipeline,
summarized in Fig.~\ref{fig:dataset_pipeline}.  DICOM pixels are decoded with
pydicom 3.0.2.  The modality LUT and, when available, the DICOM VOI LUT or
window are applied in that order; VOI processing is non-strict, so a malformed
or absent optional transform falls back to the modality-corrected pixels and is
recorded in the export metadata.  A \texttt{MONOCHROME1} image is then inverted
to the common bright-tissue-on-black convention.  Finally, values below the
per-image 0.5th percentile and above the 99.5th percentile are clipped, and the
remaining interval is mapped linearly to $[0,1]$ in float32.

We next estimate breast tissue with a robust largest-connected-component mask.
Its automatic threshold is derived from border intensity statistics and the
image dynamic range; the binary mask is morphologically opened and closed with
$7\!\times\!7$ and $21\!\times\!21$ elliptical kernels, respectively, reduced
to its largest component, and hole-filled.  The image is cropped to the mask's
bounding rectangle with 3\% padding relative to the larger breast dimension,
bounded to 32--128 pixels.  To prevent a tissue heuristic from deleting ground
truth, this rectangle and its retained mask are expanded whenever necessary to
contain every Mass box.  Pixels outside the final mask are set to zero.  Breasts
whose chest wall is detected on the right are horizontally mirrored, together
with their boxes and masks, giving a canonical chest-wall-left orientation.

This geometry follows the background evidence that removing excess air and
artifacts can increase effective lesion scale without using an unavailable
lesion-centred oracle crop \cite{ozic2023scenarios}.  In contrast to the earlier
draft, v2 also applies CLAHE: after breast cropping, masking, and mirroring, a
float32, mask-aware CLAHE transform with clip limit 2.0 and an
$8\!\times\!8$ tile grid is evaluated once on the complete processed breast.
The resulting grayscale signal is copied identically into R, G, and B before
the whole-image and window branches diverge.  Applying the same enhancement to
all branches holds the photometric recipe fixed while resolution and context
change.  No global histogram equalization, median or bilateral denoising,
geometric augmentation, or hard-negative mining is performed during image
materialization.  PNGs are quantized to 8-bit only at final encoding; the
compact PyTorch tensors retain the unquantized three-channel CHW float32 signal
in $[0,1]$.

\begin{figure}[t]
\centering
\setlength{\fboxsep}{5pt}
\fbox{\parbox{0.94\linewidth}{\centering\small
DICOM $\rightarrow$ modality/VOI processing $\rightarrow$ polarity correction
$\rightarrow$ 0.5--99.5 percentile normalization\\[2pt]
$\downarrow$\\[-1pt]
breast mask $\rightarrow$ padded breast crop $\rightarrow$ outside-mask zeroing
$\rightarrow$ canonical mirroring\\[2pt]
$\downarrow$\\[-1pt]
whole-breast float32 CLAHE $\rightarrow$ identical R/G/B channels\\[2pt]
$\downarrow$\\[-1pt]
original-size whole $+$ 1024/640 compact wholes $+$ source-resolution windows
}}
\caption{Default Research Dataset v2 pipeline.  Photometric and breast
standardization are shared before the data branch into global and local
representations.}
\label{fig:dataset_pipeline}
\end{figure}

\subsection{Multiresolution Whole Images and Virtual Windows}
\label{sec:dataset_representations}

V2 exports three annotated whole-image representations for every one of the
20,000 mammograms (Table~\ref{tab:dataset_resolution_ladder}).  The first is the
fixed-preprocessed breast at its variable, unpadded height and width.  For each
compact representation, that breast is independently placed at the top-left of
its own square canvas, with zero padding only on the right or bottom, and the
complete square is resized to either $1024\!\times\!1024$ or
$640\!\times\!640$.  OpenCV 4.10.0 area interpolation is used for downsampling
and linear interpolation if upsampling is required.  Per-image square
letterboxing preserves the breast aspect ratio and therefore removes the
geometric warping that is coupled to direct square resizing in many of the
protocols reviewed in Sec.~\ref{sec:background_resolution}.  The median
processed-breast long side is 2,728.5 pixels; the corresponding median linear
scale is 0.375 at 1,024 and 0.235 at 640.  Thus, both variants are lossy, but the
640 arm deliberately removes substantially more fine spatial evidence.  The two
compact variants are saved both as RGB PNGs and as float32 tensors; the
variable-size original is saved as an RGB PNG and provides the
source-coordinate image for virtual windows.

\begin{table}[t]
\centering
\caption{Representations in the Default Research Dataset v2 export.  Windows
are metadata-only and are extracted from the processed original at load time.}
\label{tab:dataset_resolution_ladder}
\scriptsize
\setlength{\tabcolsep}{3.2pt}
\begin{tabularx}{\columnwidth}{@{}llXl@{}}
\toprule
Representation & Output & Construction & Storage \\
\midrule
Processed whole & Variable H$\times$W & No pad or resize & RGB PNG \\
Compact whole & $1024^2$ & Per-image square, resize & PNG + float32 \\
Compact whole & $640^2$ & Per-image square, resize & PNG + float32 \\
Virtual windows & $1024^2$ & Strides 128/256/512 & Metadata \\
Virtual windows & $640^2$ & Stride 160 & Metadata \\
\bottomrule
\end{tabularx}
\end{table}

The virtual-window families retain source-resolution processed pixels rather
than rescaling the entire breast.  Window origins begin at zero and advance on
a regular grid; a window extending beyond the right or bottom edge is completed
with zeros rather than shifted back to the boundary.  The default families are
1,024-pixel windows at strides 128, 256, and 512, and 640-pixel windows at
stride 160.  Because only their locations, padding, and labels are stored, the
windows can be extracted on demand without duplicating highly overlapping
pixels.  This design exposes the scale--context trade-off discussed in
Sec.~\ref{sec:background_resolution}: compact wholes retain global anatomy but
shrink lesions, whereas windows retain local sampling density at the cost of
breast-level context.

For experiments using the optional window manifests, training retains every
eligible Mass-positive window and samples clean negative windows from breasts
with no Mass in either view toward a 1:1 positive/negative ratio using seed 123.
Validation and test retain every eligible grid window and are not balanced.  A
geometric source-extent filter rejects windows with less than 10\% in-bounds
source area in training or 5\% in validation/test; it is an edge-padding proxy,
not a second pixel-derived breast mask.  A positive window bypasses this filter
so a labeled lesion is not lost.

Mass boxes undergo exactly the same crop and mirror operations as their source
image.  Whole-image boxes remain in the fixed-preprocessed coordinate system.
For a compact whole, a coordinate is first translated by its letterbox offset
and then multiplied by that image's recorded square-to-output scale.  For a
virtual window, a box is intersected with the window and translated to local
coordinates; it is retained when at least 5\% of its original area remains
visible.  The exported manifest records every crop origin, padding amount,
mirror flag, scale, and source annotation identifier.  Whole-image annotations
are written in normalized YOLO and COCO form; virtual-window boxes are stored
in a crop-keyed annotation table for direct loading.  Dataset construction
performs no cross-window or cross-view box fusion; any such model-side operation
and its inverse mapping are specified in Sec.~\ref{sec:methods}.

\subsection{Dataset Limitations}
\label{sec:dataset_limitations}

Four limitations shape the interpretation of the experiments.  First, the
targets reflect radiologist consensus without pathology, so performance
measures agreement with localized findings rather than cancer detection.
Second, BI-RADS 2 findings are not boxed; some apparent false positives may
therefore correspond to benign findings outside the annotation policy.  Third,
per-image percentile normalization and CLAHE remove absolute intensity
comparability, while the heuristic tissue mask, crop, and canonical mirroring
change the acquisition presentation.  Mass-preserving crop expansion protects
the labeled targets, but these remain controlled engineering choices rather
than claims of optimal clinical display processing.  Finally, all examinations
come from two hospitals in one country, and the vendor mix, density
distribution, mixture of screening and diagnostic cases, and clinical reading
protocol may not represent another screening program.  The official held-out
test set supports reproducible within-dataset comparison, but external
validation is still required before making claims about cross-site or clinical
generalization.
